
\documentclass[10pt]{beamer}

\usepackage{graphicx}
\usepackage{xcolor}
\usepackage{amsfonts}
\usepackage{amsmath}
\usepackage{amssymb}
\usepackage{centernot}
\usepackage{verbatim}
\usepackage[skip=10pt plus1pt] {parskip}

%Information to be included in the title page:
\title{Formal Logic}
\author{UChicago Political Science Math Camp}
\institute{Ryan Gibbons \& Luke Cain}
\date{2025}

\begin{document}

\begin{frame}
\begin{center}
    \includegraphics[width=10cm]{/Users/ryangibbons/Downloads/IMG_4235.jpeg}
\end{center}
\end{frame}

\frame{\titlepage}


\begin{frame}{A Justification of Social Science}
    \begin{itemize}
        \item Social science as an \emph{explanatory pursuit}
        \item Evidentiary collection and analysis provides a \emph{means} to positive explanations
        \item Unexplained phenomena are unanswered questions
    \end{itemize}
\vspace{.7cm}
A central motivation: we use empirical information to produce theoretical insights, and we want our theoretical insights to be \emph{maximally applicable} while retaining accuracy
\end{frame}

\begin{frame}{Why Proofs?}
    It is almost always easier to definitively observe a phenomenon than it is to definitively explain it.

    If we cannot definitively explain it, we have not answered the important question (normative argument).

    If we cannot definitively explain it, we cannot rule out other potential mechanisms (logical argument).
\end{frame}

\begin{frame}
    \begin{center}
 \includegraphics[width=10cm]{/Users/ryangibbons/Downloads/theorem.png}
    \end{center}
\end{frame}


\begin{frame}{The Structure of Formal Logic}
    Formal logic largely focuses on three translations:
    \begin{enumerate}
        \item Empirical observations into informational inputs
        \item Informational inputs into theoretical mechanisms
        \item Theoretical mechanisms into applied inference and/or prediction
    \end{enumerate}

    We want to maximize what we can gleam from each of these three translations. 
\end{frame}


\begin{frame}{Foundational Notation}
    Formal logic requires formal notation, even beyond our set-building language. A brief overview:

    \begin{itemize}
        \item $\implies$ \quad ``implies''
        \item $\iff$ \quad ``if and only if'' \textcolor{red}{(iff)}
        \item $\neg$ \quad ``not''
        \item $\centernot\implies$ \quad ``does not imply.'' (\textbackslash{} generally means ``not.'')
    \end{itemize}

    You might occasionally see $\land$ for ``and,'' $\lor$ for ``or.'' Just write the damn words
\end{frame}

\begin{frame}{An Aside For This Lecture}
    In this lecture, I will refer to ``A,B,C...'' as \emph{individual claims}

    I will refer to \emph{logical conditions} using greek letters -- $\alpha,\beta$, etc.

    \emph{We don't apply any substantive meaning (yet).}
\end{frame}

\begin{frame}[fragile]{Implies}
Our strongest operator is $\implies$, typed as \verb|\implies|:

If $A \implies B$, then claim $A$'s truth lends claim $B$'s truth.
\pause

For example:

\begin{itemize}
    \item \verb|[state]| follows the ``Nordic Model'' $\implies$ \verb|[state]| employs a social democratic welfare system
    \item \verb|[state]| is a liberal democracy $\implies$ \verb|[state]| has a legislature \verb|OR| \verb|[state]| is a direct democracy
    \item United Nations recognizing \verb|[state]| $\implies$ \verb|[state]| is a country
\end{itemize}
\end{frame}

\begin{frame}[fragile]{Limits of Implication}
Implication is a \textcolor{red}{unidirectional} relationship: $A \implies B$ does not tell us that $B \implies A$.

Consider the claim $A \implies B$. Then, we have:
\begin{itemize}
    \item If $A$, then $B$;
    \item If $\neg B$, then $\neg A$
\end{itemize}
But we do \textcolor{red}{NOT necessarily} have:
\begin{itemize}
    \item If $B$, then $A$;
    \item If $\neg A$, then $\neg B$
\end{itemize}
\textcolor{lightgray}{What do these conditions look like formalized?}
\end{frame}

\begin{frame}{Handling Negation (pt. 1)}
    We use two different ``negation'' operators, with the following loose applications...

    \begin{itemize}
        \item We use $\not$ \quad over an operation to show that the relationship is \textbf{not necessarily} true
        \item We use $\neg$ to show that some relationship is \textbf{necessarily not} true
    \end{itemize}

    For example, the following two condition are \textcolor{red}{not} equivalent:

    \begin{enumerate}
        \item $A \centernot\implies B$
        \item $A \implies \neg B$
    \end{enumerate}
\end{frame}

\begin{frame}{Iff}
    Our strongest operator is $\iff$, or ``if and only if.'' This implies a \textbf{bidirectional relationship}.

    If $A \iff B$, then:
    \begin{enumerate}
        \item $A \implies B$, AND 
        \item $B \implies A$
    \end{enumerate}

    Neither $A$ nor $B$ can be true on their own. 
    
    \emph{If we know $A \iff B$ to be true, what do we need to observe to determine either $A$ or $B$?}
\end{frame}

\begin{frame}[fragile]{Iff Examples}
    Iff statements are especially present when defining terms and concepts. \emph{\textcolor{lightgray}{Why?}} 

    \begin{itemize}
        \item \verb|[state]| has free \& fair elections with uncertain outcomes $\iff$ \verb|[state]| is a democracy
        \item A central bank is entirely independent $\iff$ elected officials cannot control or influence the interest rate
        \item A legislative bloc has an outright majority $\iff$ it holds more than half of legislative seats
    \end{itemize}
\end{frame}


\begin{frame}{Logical Grammar}
    Logical statements often require complexity. We use parentheses to identify an order of operations:
    \begin{itemize}
        \item $(A \text{ and } B) \implies C$: If $A$ and $B$, then $C$.
        \item $A \implies (B \text{ and } C)$: If $A$, then $B$ and $C$.
    \end{itemize}
    
    This looks simple so far...
\end{frame}

\begin{frame}{Handling Negation (pt. 2)}
    Logical ordering becomes very important when considering negative conditions or interactions. For example, the following statements are \textcolor{red}{not} equivalent:

    \begin{enumerate}
        \item $\neg (A \text { and } B) \implies C$
        \item $\neg A \text { and } \neg B \implies C$
    \end{enumerate}

    \emph{Which of these statements would be more challenging to prove?}
\end{frame}


\begin{frame}{Strength of Conditions}
    Some conditions are harder to meet than others. A helpful consideration is identifying \textbf{nesting} conditions: is one condition implies by another?

    We can consider a condition $\alpha$ as \textbf{stronger} than a condition $\beta$ if \textcolor{red}{(not iff)} $\alpha \implies \beta$.

    \pause
    For example, the following conditions have increasing strength:
    \begin{enumerate}
        \item A legislative bloc has plurality control
        \item A legislative bloc is the dominant group of a majority coalition
        \item A legislative bloc has an outright majority
        \item A legislative bloc has a supermajority (can pass any item without failure)
    \end{enumerate}
\end{frame}

\begin{frame}{An Interesting Application: Bureaucratic Discretion}
    A common question in domestic politics regards how much \emph{discretion} principals lend to their agents (see Huber \& Shipman 2002).

    Central tradeoff: providing discretion allows for both effectiveness and shirking.

    But, to evaluate this theory, we must measure bureaucratic discretion...
\end{frame}

\begin{frame}{An Interesting Application: Bureaucratic Discretion}
    Much of the literature on bureaucratic discretion uses legislative word count as an independent variable (again, see Huber \& Shipman 2002). This is a proxy in accordance with the following logic:

    \begin{itemize}
        \item A bureaucrat facing conditions $A$ and $B$ is more constrained than a bureaucrat facing only $A$ or only $B$
        \item In order to implement conditions $A$ and $B$, legislators often need to write \emph{at least} $A$ \emph{plus} some additions to achieve $B$
        \item Therefore, in expectation, longer written laws should imply additional conditions on bureaucratic behavior (and therefore less discretion)
    \end{itemize}
\end{frame}

\begin{frame}{Back to It}
    Clearly, $\iff$ is a stronger condition than $\implies$. (Why?)

    \pause
    Is $\centernot\iff$ a stronger condition than $\iff$?

    \pause
    Is $\centernot\iff$ a stronger condition than $\centernot\implies$?

    \pause
    Is $\neg \iff \neg$ a stronger condition than $\centernot\implies$?
\end{frame}


\begin{frame}{Nesting: Necessary \& Sufficient}
    We can characterize conditions as \textbf{necessary} and/or \textbf{sufficient}:

    \begin{itemize}
        \item Condition $\alpha$ is sufficient for condition $\beta$ if $\alpha \implies \beta$.
        \begin{itemize}
            \item \textcolor{gray}{If this is true, what does observing $\not \beta$ tell us?}
        \end{itemize}
        \item Condition $\alpha$ is necessary for condition $\beta$ if $\neg \alpha \implies \neg \beta$
        \begin{itemize}
            \item \textcolor{gray}{If this is true, what does observing $\beta$ tell us?}
        \end{itemize}
    \end{itemize}

    \emph{Which type of condition is stronger?}
\end{frame}



\begin{frame}{What Is To Be Done?}
    To produce meaningful conclusions, we will focus on four fundamental ``proofs'' using these logical operations:

    \begin{enumerate}
        \item Direct
        \item Contrapositive
        \item Contradiction
        \item Cases
    \end{enumerate}

    These are general strategies; there's not always a ``correct'' fit.
\end{frame}


\begin{frame}{Direct Proofs}
    Direct proofs simply show that multiple conditions are equivalent.

    Suppose we want to prove $A \iff B$. Then, a direct proof would need to show both of the following:

    \begin{itemize}
        \item If $A$, then $B$ ($A \implies B$)
        \item If $B$, then $A$ ($B \implies A$)
    \end{itemize}
\end{frame}

\begin{frame}{Why Not Direct?}
    Direct proofs can be great \emph{if we have sufficient information}; we don't always.

    For example, suppose we want to prove that every perfect direct democracy\footnote{Every citizen has equal influence over outcomes} has an imbalanced tiebreak institution. Can we prove this directly?

    %answer: no, because whether or not a nation is a perfect direct democracy presumably depends on whether or not the tiebreak is objective
\end{frame}

\begin{frame}{The Contrapositive}
    Contrapositives invert $\implies$ and allows us to prove via \textbf{falsification}. Namely, $A \implies B$ is equivalent to $\neg B \implies \neg A$.

    Suppose we want to prove $A \iff B$, but could not conclude directly that $B \implies A$ Then, a contrapositive proof would need to show:

    \begin{itemize}
        \item If $A$, then $B$ (or $A \implies B$)
        \item If $\neg B$, then $\neg A$ (or $\neg B \implies \neg A$)
    \end{itemize}

    This tends to be an easier proof of $\iff$ statements than the direct route above.
\end{frame}

\begin{frame}{A Contrapositive Example}
    Can we use a contrapositive proof to prove our above claim (that every perfect direct democracy has an imbalanced tiebreak institution)?
    \pause

    Yes! We just need to show the following: If a state does NOT have an imbalanced tiebreak institution, then it is NOT a perfect direct democracy. We can prove this directly.
\end{frame}


\begin{frame}{Proof by Contradiction}
    Contradiction is a powerful proof technique to demonstrate that a claim is true \emph{in all cases.}

    We start with a claim that contradicts what we want to prove. If we want to prove $A \iff B$, we then start with: ``Suppose not. Then...''

    Think of contradiction as a rabbit-hole proof: expand conditions as far as we can, \underline{until something breaks}

\end{frame}


\begin{frame}{Contradiction in Abstract}
    Let's expand upon the arbitrary proof by contradiction (proving that $A \iff B$). The proof might proceed:

    \begin{enumerate}
        \item Suppose not.
        \item Then, either $A \centernot\implies B$ or $B \centernot\implies A$.
        \item Then, either there must exist a case satisfying only $A$ and not $B$ or only $B$ and not $A$.
        \item \emph{Disprove either claim.}
    \end{enumerate}
\end{frame}

\begin{frame}{Using Contradiction: MVT}
    The Median Voter Theorem suggests that candidate platforms will converge to the preference of the median voter:

    \begin{quote}
        Suppose there are two candidates $A$ and $B$ who establish policy positions along a single dimension (canonically, a preferred tax rate). Then, if $A$ and $B$ are only motivated by winning office, they will both establish policy positions equivalent to the preference of the mathematically median voter.
    \end{quote}
\end{frame}


\begin{frame}{Why All The Excitement?}
    These proof structures have gotten increasingly more complex -- but, they have widened the range of relationships we are examining, and we are usually just looking for one to ``break.''

    From $A \iff B$, we've extended to include:
    \begin{itemize}
        \item Direct: $A \implies B$, $B \implies A$
        \item Contrapositive: $A \implies B$, $\neg B \implies \neg A$
        \item Contradiction: $A \centernot\implies B$, $B \centernot\implies A$
    \end{itemize}

    Proof structures allow us to cast wider nets when fishing for claims.
\end{frame}


\begin{frame}{Proof by Cases}
Not necessarily mutually exclusive from the above, but an important consideration: oftentimes, we may only have a finite number of possible cases, and can use that to our advantage.

Suppose we want to prove that every party represented in Congress holds a primary. We have two options:

\begin{enumerate}
    \item Prove that in order to be represented in Congress, a party must have held a primary
    \item Just look at the parties represented in Congress (2ish), and see whether or not they held a primary
\end{enumerate}
\end{frame}

\begin{frame}[fragile]{Proof By Cases}
    One common case structure we can often use as a shortcut is derived from ordering. Suppose we have two parties, \verb|LUKE| and \verb|RYAN|. Then, one of the following MUST be true w.r.t. vote counts:
    \begin{enumerate}
        \item \verb|#LUKE| $>$ \verb|#RYAN|
        \item \verb|#LUKE| $<$ \verb|#RYAN| (most likely)
        \item \verb|#LUKE| = \verb|#RYAN|
    \end{enumerate}
    
    \emph{Not every case needs to be proven the same way...}
\end{frame}


\begin{frame}[fragile]{Axiomatic Truths}

    Our goal is usually to prove ``as far down as possible'' -- but, we need some initial foundations.

    \textbf{Axioms} are things we hold true without proofs; they are generally extremely basic.
    
\end{frame}



\end{document}